\section{Participação em Eventos Científicos e Atividades Acadêmicas}
\label{sec:atividades-academicas}

\subsection{Congresso Acadêmico de 2026}

Como atividade de divulgação científica, uma versão preliminar da pesquisa foi apresentada no Congresso Acadêmico em 14 de agosto de 2026. A apresentação contextualizou a base de alarmes industriais, as análises temporais e a definição do alvo como o volume diário de recorrências com intervalo (IAT) inferior a 30 minutos. Na ocasião, expuseram-se experimentos iniciais (com 17 e 23 preditores e Random Forest) cujas inconsistências motivaram a reformulação do protocolo experimental, do tratamento de hiatos e da validação temporal.  Posteriormente, a modelagem foi consolidada comparando conjuntos de 16 (simples) e 36 preditores (completo) via Random Forest e XGBoost sob validação expansiva e teste isolado. O XGBoost completo foi selecionado como modelo final (MAE de 8,817 e $R^2$ de 0,051), enquanto o XGBoost simples (MAE de 9,267) despontou como alternativa viável para implantação industrial por dispensar dados externos e de máquinas individuais. Apesar do baixo erro médio, a capacidade explicativa limitada restringiu o uso do modelo ao apoio e priorização do monitoramento operacional, sem prever picos extremos com precisão.  Durante a arguição, esclareceu-se à banca que a série diária foi construída tratando dias zerados explicitamente como potenciais indisponibilidades, e que as predições utilizam apenas informações anteriores ao dia-alvo. Com essas justificativas, a banca considerou o projeto satisfatório.  
